% SOURCE_STATEMENT: "Find general conditions under which scaling-limit and tensor-product operations commute for nets in algebraic quantum field theory."
% Preparation and final self-audit: CONTRIBUTING.md prompts C and D; source/audit artifacts retained in the local physics_statements_2026_10_11 batch.
% !TeX program = lualatex
\documentclass[11pt]{article}
\usepackage[margin=25mm]{geometry}
\usepackage{fontspec}
\setmainfont{Latin Modern Roman}
\usepackage{amsmath,amssymb,mathtools}
\usepackage{unicode-math}
\setmathfont{Latin Modern Math}
\usepackage{xurl,hyperref}
\hypersetup{colorlinks=true,urlcolor=blue}
\setcounter{secnumdepth}{0}
\setlength{\parindent}{0pt}
\setlength{\parskip}{5pt}
\setlength{\emergencystretch}{3em}
\newcommand{\sref}[2]{\hyperlink{source:#1:#2}{[S#2]}}
\newcommand{\cataloguescope}[1]{\paragraph{Recorded target.}#1}
\newcommand{\R}{\mathbb R}
\newcommand{\C}{\mathbb C}
\newcommand{\N}{\mathbb N}
\newcommand{\Z}{\mathbb Z}
\newcommand{\Q}{\mathbb Q}
\newcommand{\F}{\mathbb F}
\newcommand{\eps}{\varepsilon}
\begin{document}
\section{Scaling Limits of Tensor Products in AQFT}\label{statement}

\subsection{Definitions and mathematical statement}

\paragraph{Shared notation.}$\mathbb N=\{1,2,\ldots\}$; the standard number fields have their usual meanings. All additional conventions are stated locally.


\paragraph{Definitions and conventions.} An algebraic quantum field theory net assigns an algebra $\mathcal F(O)$ to each spacetime region, together with vacuum and translation data. The tensor-product net has local algebras $\mathcal F_1(O)\overline\otimes\mathcal F_2(O)$. A scaling limit is obtained by rescaling regions $\lambda O$ and taking a scaling-algebra state limit as $\lambda\downarrow0$. Commutation means that the corresponding limit net is isomorphic to $\mathcal F_{1,0}\overline\otimes\mathcal F_{2,0}$, with compatible choices of limit states.

\paragraph{Source scope and qualifications.} The report proves commutation under a uniform small-p nuclearity bound and twisted Haag duality. The open request is for more general conditions, not to reprove that stated sufficient case or to assert commutation without phase-space assumptions.

\paragraph{Catalogue mathematical target.}
\paragraph{Catalogue wording retained for traceability.}
Find general conditions under which scaling-limit and tensor-product operations commute for nets in algebraic quantum field theory.

\paragraph{Saved literature context.} The supplied catalogue's assessment dated 2026-08-21 records: The report explicitly asks for general conditions allowing scaling limits and tensor products to interchange; the extracted question remains too general for a verified resolution. The assessment is retained as dated provenance; it is not a new comprehensive literature or open-status certification.
\subsection{Short English statement}
Find broader hypotheses under which taking the ultraviolet limit of two quantum field theories commutes with forming their tensor product.
\subsection{Sources}
\begin{itemize}
\item[\hypertarget{source:statement:1}{S1}] ULAM.AI and the UnsolvedMath Contributors, UnsolvedMath, catalogue entry. Catalogue curation under CC BY 4.0; definitions and English exposition added locally.
\url{https://huggingface.co/datasets/ulamai/UnsolvedMath}.
\item[\hypertarget{source:statement:2}{S2}] Noncommutative Geometry and Quantum Field Theory (2005), Oberwolfach Reports, DOI 10.4171/owr/2005/48, pp. 2705--2760; official report source and mathematical context.
\url{https://doi.org/10.4171/owr/2005/48}.
\item[\hypertarget{source:statement:3}{S3}] Official EMS Press PDF of the same Oberwolfach report; PDF pages 50--51, printed pages 2754--2755.
\url{https://ems.press/content/serial-article-files/46022}.
\end{itemize}
\label{end:statement}
\end{document}
